
Professional medicine questions appear in The Pile v1 at approximately 43 times the rate of high school mathematics questions — yet standard practice reports MMLU accuracy as a single number averaged across all 57 sub-tasks, as though contamination were uniform. The experimental results reported here demonstrate that it is not.

Training data contamination, defined here as the n-gram overlap between benchmark test sets and training corpora, is acknowledged as a potential source of performance inflation in large language model evaluation. A model trained on text that contains benchmark questions has, in effect, been exposed to the answers. The NLP community has responded with contamination reporting norms, decontamination filtering, and tooling such as WIMBD \citep{Elazaretal2023}. These responses share a common limitation: they treat contamination as a scalar per benchmark rather than a structured, sub-task-level phenomenon.

This framing may obscure a deeper pattern. The three most widely used open training corpora — The Pile v1 \citep{Gaoetal2020}, C4 en.noclean \citep{Dodgeetal2021}, and RedPajama-v1 [TogetherComputer, 2023] — differ substantially in source composition. The Pile deliberately includes domain-aligned academic sources (PubMed Central, ArXiv, FreeLaw). C4 applies quality filtering to Common Crawl that incidentally removes content overlapping with NLP benchmarks. RedPajama combines web text, code, academic text, and books. If source composition shapes which benchmark sub-tasks are most contaminated, then aggregate contamination rates obscure exactly the information practitioners need to assess evaluation validity.

No published work provides a unified cross-corpus contamination matrix — contamination rates for each benchmark sub-task against each of the three major open training corpora, computed with consistent methodology. WIMBD \citep{Elazaretal2023} provides 13-gram containment tooling for The Pile but does not extend to C4 or RedPajama. The GPT-4 Technical Report [OpenAI, 2023] reports contamination using Jaccard similarity against a closed corpus. Dodge et al. [2021] document C4's properties but do not measure contamination against MMLU, HellaSwag, or BIG-Bench Hard. The contamination atlas — a structured map of which sub-tasks carry elevated contamination risk from which corpora — does not exist in prior work.

The central hypothesis of this study is that corpus source composition, not merely corpus scale, is associated with which benchmark domains are most contaminated. Quality-filtered web text (C4) shows lower mean contamination than The Pile despite similar scale. The Pile and RedPajama show statistically indistinguishable contamination distributions because both draw substantially from CommonCrawl web text. This hypothesis is tested empirically via a 59-sub-task × 3-corpus contamination matrix — 177 (sub-task, corpus) cells — using 13-gram containment methodology validated against WIMBD published baselines (Spearman ρ=0.721).

The contributions of this work are:

\begin{enumerate}
\item \textbf{The Cross-Corpus Contamination Matrix}: A 59-sub-task × 3-corpus 13-gram contamination matrix for MMLU, HellaSwag, and BIG-Bench Hard against The Pile v1, C4 en.noclean, and RedPajama-v1, computed with a consistent, open-source methodology and validated against WIMBD published rates.
\end{enumerate}

\begin{enumerate}
\item \textbf{The Corpus Composition Effect}: Empirical demonstration that corpus source composition is associated with contamination levels. C4's quality filtering is associated with a 38\% lower mean contamination rate relative to The Pile (directional ordering robust to ρ>0.995 sensitivity analysis; absolute values carry ~10–20\% uncertainty from literature-calibrated scaling factors applied to 10\% corpus samples). The Pile-RedPajama statistical equivalence (p=0.810) is a novel finding consistent with shared CommonCrawl ancestry.
\end{enumerate}

\begin{enumerate}
\item \textbf{Domain-Stratified Contamination Profiles} (exploratory): Evidence that academic/professional MMLU sub-tasks show approximately 1.9× higher contamination rates than commonsense benchmarks across all three corpora (Cohen's d=0.85). This finding is supported by a Kruskal-Wallis interaction test (H=22.08, p=0.0005) but was not confirmed by the pre-registered directional Mann-Whitney test due to insufficient commonsense category granularity (n=2 sub-tasks). This result should be considered exploratory.
\end{enumerate}

\begin{enumerate}
\item \textbf{Methodological Validation}: Demonstration that rank-based contamination analysis is robust to text format variations (Spearman ρ=0.74) and corpus sampling fractions (ρ>0.995).
\end{enumerate}

